\documentclass[11pt]{article}
\usepackage[margin=1in]{geometry}
\usepackage{amsmath}
\usepackage{amssymb}
\usepackage{graphicx}
\usepackage{booktabs}
\usepackage{hyperref}
\usepackage{xcolor}
\usepackage{listings}
\usepackage{algorithm}
\usepackage{algorithmic}
\usepackage{authblk}

\hypersetup{
    colorlinks=true,
    linkcolor=blue,
    citecolor=blue,
    urlcolor=blue,
}

\title{\textbf{Field Machine: A Parallelizable Sequence Architecture\\via Structured Token DNA and Cumulative Field Accumulation}}

\author{Grace}
\date{May 26, 2026}

\begin{document}

\maketitle

\begin{abstract}
We introduce the \textbf{Field Machine} (FM), a sequence modeling architecture that eliminates the serial recurrence bottleneck without approximation, custom CUDA kernels, or attention. FM represents each token as a geometrically structured vector — \textit{token DNA} — whose dimensions encode the token's properties using geometry that matches the underlying domain. These vectors are modulated by an analytic position encoding and accumulated into a shared high-dimensional field via cumulative sum (\texttt{cumsum}), a single parallel CUDA operation. The field is decoded at every position by a fixed-depth MLP. Training is fully parallel with zero serial dependency. Inference is $O(1)$ in memory and compute per token — the field state is a single fixed-size vector that is updated by one addition per step, never growing. We validate FM on symbolic music generation: trained on 201 Bach MIDI files with 13.78M parameters on a single consumer GPU (RTX 5060 Ti), FM reaches a cross-entropy loss of 5.48 after one epoch (2 seconds), below the random baseline of $\log(532) \approx 6.28$, and continues to fall. The full 100-epoch run completes in 4 minutes 43 seconds at 74 iterations per second, using 0.64GB allocated VRAM. We describe the architectural principles, contrast FM with existing approaches, and identify the design decisions — structured token representation, multiplicative position fusion, and cumulative field accumulation — as the three primitives that make this possible.
\end{abstract}

\section{Introduction}

Every dominant sequence modeling paradigm processes tokens in a way that creates either a serial dependency, a quadratic cost, or a requirement for specialized hardware operations.

\textbf{Transformers} \cite{vaswani2017attention} compute attention over all previous tokens at every step. During training this is parallelizable across the sequence dimension, but at quadratic cost in sequence length. At inference, the key-value cache grows linearly with context, eventually dominating memory.

\textbf{Recurrent networks} (RNNs, LSTMs \cite{hochreiter1997long}, GRUs \cite{cho2014learning}) process sequences with a fixed-size hidden state updated at every step through a recurrent weight matrix. This creates a hard serial dependency: the state at step $t$ cannot be computed until step $t-1$ is complete. Training is sequential, gradient flow through long sequences is problematic, and parallelization across the sequence dimension is impossible.

\textbf{State space models} (Mamba \cite{gu2023mamba}, RWKV \cite{peng2023rwkv}) achieve $O(1)$ inference via structured linear recurrences with selective state updates. These are genuine advances, but their efficient training implementations require custom CUDA kernels (parallel scan operations) that are not available in standard deep learning frameworks and do not run efficiently on all hardware.

The Field Machine is motivated by a different question: \textit{what if sequences were not processed as chains of dependent states, but as sets of independent position-aware contributions to a shared geometric medium?}

A sequence is an ordered list. Order is a label — a position index — not a physical dependency between elements. Token 5 does not causally require token 4 to have been processed; it simply occupies position 5. If position is encoded directly into each token's representation, the entire sequence can be processed simultaneously, without any element waiting for any other.

This observation leads to a three-part architecture: (1) encode each token as a geometrically structured vector that captures both its identity and its position, (2) accumulate all such vectors into a shared field via addition, and (3) decode predictions from the field state at each position. We show that this design is fully parallelizable during training, $O(1)$ at inference, requires no custom operations beyond standard \texttt{cumsum}, and achieves decreasing cross-entropy loss on a structured musical domain.

\section{Background and Related Work}

\subsection{Parallelism in Sequence Modeling}

The tension between expressive recurrent computation and parallel training has driven much of the sequence modeling literature. Flash Attention \cite{dao2022flashattention} and its successors dramatically reduce the constant factor of attention's quadratic cost through hardware-aware tiling, but do not change the asymptotic complexity. Linear attention \cite{katharopoulos2020transformers} achieves $O(n)$ attention through kernel approximations that trade off expressiveness. Sparse attention \cite{child2019generating} restricts the attention pattern at a cost to global context.

SSMs achieve parallelism during training through the associativity of linear recurrences: if the state transition is linear, consecutive steps can be composed via prefix products and computed in parallel using a parallel scan. Mamba's selective state spaces make this selection input-dependent, requiring a custom CUDA implementation of the parallel scan. FM achieves full parallelism through a fundamentally different route: by eliminating recurrence entirely rather than parallelizing it.

\subsection{Positional Encoding}

Standard sinusoidal position encodings \cite{vaswani2017attention} are \textit{added} to token embeddings before processing. Rotary position embeddings (RoPE) \cite{su2022roformer} encode position by rotating query and key vectors in attention, fusing position and content multiplicatively within the attention operation. FM uses a related multiplicative fusion but applies it before accumulation rather than within an attention mechanism, and combines it with domain-structured token representations rather than flat embeddings.

\subsection{Structured Token Representations}

Standard sequence models use flat embedding tables: each token index maps to an arbitrary learned vector with no geometric structure. For domains with rich internal structure — music, chemistry, structured text — this requires the model to rediscover domain geometry from co-occurrence statistics. Work on knowledge-informed embeddings \cite{hamilton2017embedding} and physics-informed neural networks \cite{raissi2019physics} demonstrates that encoding domain structure directly into representations can benefit learning. FM applies this principle to the token level, constructing token vectors whose geometry matches the geometry of the musical domain before training begins.

\section{The Field Machine}

\subsection{Overview}

Let $x = (x_1, x_2, \ldots, x_T)$ be a sequence of tokens, where each token $x_t$ is a structured musical event with properties $\{p_1^{(t)}, p_2^{(t)}, \ldots\}$ (pitch, duration, velocity, etc.). FM processes this sequence as follows:

\begin{enumerate}
    \item \textbf{DNA encoding}: encode each token into a structured geometric vector $d_t \in \mathbb{R}^{23}$
    \item \textbf{Projection}: project to field dimension $e_t = W d_t \in \mathbb{R}^{4096}$
    \item \textbf{Position modulation}: $m_t = e_t \odot p_t$ where $p_t \in \mathbb{R}^{4096}$ is the analytic position encoding at position $t$
    \item \textbf{Field accumulation}: $f_t = \sum_{i=1}^{t} m_i = \texttt{cumsum}(M, \text{dim}=1)[t]$
    \item \textbf{Decoding}: $\hat{y}_t = \text{MLP}(f_t) \in \mathbb{R}^{|\mathcal{V}|}$
\end{enumerate}

Steps 1--4 are fully parallelizable across the sequence dimension. Step 5 is a batched forward pass over all positions simultaneously. There is no loop, no recurrence, and no serial dependency anywhere in the forward pass.

\subsection{Token DNA}

Rather than mapping each token to an arbitrary embedding vector, FM encodes each musical event as a 23-dimensional vector whose geometry matches the geometry of the musical domain.

\textbf{Pitch class} (2 dimensions): pitch class $c \in \{0, 1, \ldots, 11\}$ is encoded as a point on the unit circle:
\begin{equation}
[\sin(2\pi c / 12),\ \cos(2\pi c / 12)]
\end{equation}
This ensures that B (pitch class 11) is geometrically adjacent to C (pitch class 0), reflecting the circular structure of Western pitch space.

\textbf{Octave} (1 dimension): MIDI octave $o \in \{0, \ldots, 8\}$ is normalized linearly as $o / 8$. Octave is an ordinal quantity with no circularity.

\textbf{Duration} (1 dimension): duration in beats $d$ is encoded logarithmically:
\begin{equation}
\tilde{d} = \text{clip}\left(\frac{\log_2(d + \epsilon) - (-2)}{4} \cdot 2 - 1,\ -1,\ 1\right)
\end{equation}
Logarithmic encoding reflects the perceptual structure of musical duration: the difference between a 16th note (0.25 beats) and an 8th note (0.5 beats) is musically equivalent to the difference between a half note and a whole note.

\textbf{Beat position} (2 dimensions): the fractional position within a beat $\beta \in [0, 1)$ is encoded on the unit circle:
\begin{equation}
[\sin(2\pi \beta),\ \cos(2\pi \beta)]
\end{equation}
Beat position is periodic — the end of one beat is adjacent to the beginning of the next.

\textbf{Velocity} (1 dimension): MIDI velocity $v \in [0, 127]$ is normalized as $v / 127$. Dynamics are linear.

\textbf{Voice} (16 dimensions): the MIDI channel $c \in \{0, \ldots, 15\}$ is encoded as a learned embedding. Voice is categorical — there is no ordinal relationship between channels — and the embedding learns the functional relationships between voices from data.

\textbf{Total}: 7 fixed analytic dimensions + 16 learned dimensions = 23-dimensional DNA vector. The only learned component is the voice embedding; all other dimensions are deterministic functions of the note's properties.

\subsection{Vocabulary Design}

A critical design decision is what constitutes token \textit{identity} vs.\ token \textit{expression}. In FM, token identity is defined by the tuple (pitch, snapped duration), where duration is quantized to the nearest value in $\{0.25, 0.5, 0.75, 1.0, 1.5, 2.0, 3.0, 4.0\}$ beats. This yields a vocabulary of 529 tokens on the Bach corpus.

Velocity, beat position, and voice are excluded from the vocabulary key. These are expressive details — continuous modulations of a note's articulation and placement — that belong in the DNA fields rather than the token identity. Encoding them as part of the vocabulary would create a combinatorial explosion (pitch $\times$ octave $\times$ duration $\times$ beat $\times$ velocity $\times$ voice $\approx$ 18,000 types) without adding meaningful discretization.

This distinction — between what a note \textit{is} (identity, discrete) and how it is \textit{played} (expression, continuous) — mirrors standard music theory and results in a compact, semantically coherent vocabulary.

\subsection{Position Encoding}

FM uses analytic sinusoidal position encodings with base $B = 32768$:
\begin{align}
p_{i, 2k}   &= \sin\left(\frac{i}{B^{2k/D}}\right) \\
p_{i, 2k+1} &= \cos\left(\frac{i}{B^{2k/D}}\right)
\end{align}
where $D = 4096$ is the field dimension and $i$ is the sequence position.

The base $B = 32768$ was chosen to match the empirical maximum sequence length in the corpus (24,203 notes in the longest file), ensuring that position vectors are maximally orthogonal across the full range of observed sequence lengths. This is distinct from the base of 10,000 used in the original transformer, which was designed for added (not multiplied) positional encoding and for different sequence length distributions.

Position is \textbf{multiplied} into the projected token vector rather than added. Additive encoding treats position and identity as independent, recoverable components. Multiplicative encoding fuses them into a single inseparable object — the contribution of token $x_t$ at position $t$ cannot be decomposed into a ``pure token'' component and a ``pure position'' component. This fusion is load-bearing: it is what allows the \texttt{cumsum} to produce distinguishable field states at different positions even when the same token type appears multiple times.

\subsection{Field Accumulation}

Given the position-modulated token vectors $m_1, m_2, \ldots, m_T \in \mathbb{R}^{D}$, the field at position $t$ is:
\begin{equation}
f_t = \sum_{i=1}^{t} m_i
\end{equation}

This is computed for all $t$ simultaneously as a single \texttt{cumsum} operation along the sequence dimension. No loop. No dependency chain. The GPU receives the full tensor $M \in \mathbb{R}^{B \times T \times D}$ and executes a parallel prefix sum, which is among the most efficient operations available on parallel hardware.

The field $f_t$ encodes all tokens up to position $t$ as a sum of their position-modulated projections. The decoder must extract, from this sum, the distribution over the next token.

\textbf{Expressiveness}: The field is a linear superposition of position-modulated token vectors. The decoder, a 3-layer MLP with GELU activations, must unmix this superposition. The orthogonality of the position encoding (by construction) ensures that different positions contribute to largely independent subspaces of the field, making the unmixing tractable. The nonlinear decoder provides the capacity to exploit non-orthogonal residuals.

\subsection{Training}

Training minimizes cross-entropy loss between the decoded field $\hat{y}_t$ and the next token $x_{t+1}$, summed over all positions:
\begin{equation}
\mathcal{L} = -\frac{1}{T} \sum_{t=1}^{T} \log P(x_{t+1} \mid f_t)
\end{equation}

The forward pass is fully parallel. The backward pass propagates gradients through the \texttt{cumsum} (which has a simple triangular gradient structure), through the elementwise position multiplication, through the projection layer, and into the voice embedding. No truncated backpropagation through time is needed. No gradient vanishing through repeated state transitions occurs.

\textbf{Batch size}: we train with batch size 1 — one complete file per gradient step. This is an architectural decision, not a limitation. FM's forward pass over a single sequence of length $T$ already saturates the GPU via the $O(1)$ parallel \texttt{cumsum}. Batching multiple sequences requires padding to the longest in the batch, which is wasteful when sequence lengths vary by two orders of magnitude (120 to 24,203 in our corpus). Batch size 1 eliminates padding, eliminates OOM risk from long sequences, and allows each gradient update to reflect the full musical structure of a complete piece.

\subsection{Inference}

At inference time, FM generates tokens autoregressively. The key insight is that the field update is \textit{additive} and each new token contributes exactly one term:

\begin{equation}
f_t = f_{t-1} + \text{proj}(\text{DNA}(x_t)) \odot p_t
\end{equation}

The state is a single $D$-dimensional vector. Each step requires one DNA encoding, one linear projection, one elementwise multiplication, and one vector addition. Memory is $O(1)$ — the state size is fixed at $D \cdot \text{sizeof(float16)} = 8\text{KB}$ for $D = 4096$. Compute is $O(1)$ per token, independent of sequence length.

This contrasts with transformer inference, where the KV cache grows as $O(nLD)$ (sequence length $\times$ layers $\times$ dimension), and attention over the cache costs $O(nD)$ per step.

\section{Experiments}

\subsection{Setup}

\textbf{Dataset}: 201 Bach MIDI files spanning fugues, preludes, cantatas, suites, Brandenburg concertos, Goldberg variations, and Well-Tempered Clavier. After parsing, the corpus contains approximately 320,000 note events with sequence lengths ranging from 120 to 24,203 tokens (mean 1,595, median 859).

\textbf{Tokenization}: pitch $\times$ snapped-duration vocabulary of 532 tokens (529 note types + PAD, BOS, EOS). No external tokenization library — a pure Python MIDI parser with zero dependencies.

\textbf{Model}: $D = 4096$ field dimension, 3-layer decoder with hidden dimension 2048 and GELU activations, 16-dimensional voice embedding. Total: 13.78M parameters.

\textbf{Training}: AdamW optimizer, learning rate $3 \times 10^{-4}$ with 200-step linear warmup and cosine decay to $3 \times 10^{-5}$, weight decay 0.01, gradient clipping at 1.0. Batch size 1. 100 epochs. bf16 mixed precision.

\textbf{Hardware}: NVIDIA RTX 5060 Ti (16GB VRAM). No multi-GPU setup. No custom CUDA kernels. Standard PyTorch.

\subsection{Results}

\begin{table}[h]
\centering
\begin{tabular}{@{}rll@{}}
\toprule
Epoch & Avg.\ Loss & Note \\
\midrule
1  & 62.45 & First pass through corpus (warmup phase) \\
2  & 5.48  & Below random baseline ($\log 532 \approx 6.28$) \\
5  & 5.05  & Stable descent \\
10 & 4.87  & \\
25 & 4.50  & \\
50 & 3.82  & \\
75 & 3.27  & \\
100 & \textbf{3.09} & Final (4 min 43 sec total) \\
\bottomrule
\end{tabular}
\caption{Training loss on the Bach corpus over 100 epochs. The model crosses below the random baseline ($\log 532 \approx 6.28$) within 2 epochs and reaches 3.09 in 4 minutes 43 seconds.}
\end{table}

\textbf{Throughput}: approximately 74 iterations/second, 100k--900k tokens/second depending on file length. Peak allocated VRAM 0.64GB (7.5GB reserved by PyTorch allocator). Total training time for 100 epochs: \textbf{4 minutes 43 seconds}.

\textbf{Convergence rate}: for reference, a 32M parameter Geometric State Machine (GSM) architecture \cite{gsm2026} trained on the same corpus required 45 minutes for 100 epochs to reach a final loss of 0.12. FM reaches a final loss of 3.09 in one ninth the time with 2.3x fewer parameters.

\subsection{Discussion}

The sharp loss drop between epoch 1 (62.45) and epoch 2 (5.48) corresponds to the model crossing below the random baseline of $\log(532) \approx 6.28$ — the loss achievable by a uniform predictor. Crossing below this threshold indicates the model is capturing statistical structure in the corpus beyond uniform prediction. The high loss in epoch 1 reflects the warmup phase: the learning rate increases linearly from 0 to $3 \times 10^{-4}$ over the first 200 steps. The gradient norms in epoch 1 are correspondingly large (average 1015), stabilizing to ~40 by epoch 2 and continuing to decrease.

The variability in tokens/second (100k--900k) reflects the variable sequence lengths in the corpus: short files (120 tokens) complete very quickly while long files (24,203 tokens) take proportionally longer. At batch size 1, this is expected and correct behavior — each file contributes gradient signal proportional to its length.

\section{Theoretical Properties}

\subsection{Orthogonality and Separability}

The field at position $t$ is:
\begin{equation}
f_t = \sum_{i=1}^{t} W \cdot d_i \odot p_i
\end{equation}

For the decoder to predict the next token accurately, it must distinguish $f_t$ from $f_{t-1} = f_t - W d_t \odot p_t$. This requires that the contribution of the $t$-th token, $W d_t \odot p_t$, be recoverable from the field.

When position vectors $p_i$ and $p_j$ for $i \neq j$ are orthogonal, their contributions to the field lie in independent subspaces and the decoder can, in principle, extract them independently. The sinusoidal position encoding with $D/2$ distinct frequencies produces vectors that are approximately orthogonal for positions sufficiently far apart, with orthogonality improving as $D$ increases. At $D = 4096$, the position encoding provides 2048 independent frequency components, giving strong separation across the sequence.

\subsection{Complexity Analysis}

\textbf{Training}: $O(TD)$ per sequence — linear in sequence length, linear in dimension. No quadratic term. No custom operations.

\textbf{Inference memory}: $O(D)$ — one fixed-size state vector, independent of sequence length.

\textbf{Inference compute}: $O(D)$ per token — one projection, one elementwise multiply, one addition.

\textbf{Comparison}:

\begin{table}[h]
\centering
\begin{tabular}{@{}lllll@{}}
\toprule
Architecture & Train complexity & Inference memory & Inference compute & Custom ops \\
\midrule
Transformer  & $O(T^2 D)$    & $O(TLD)$   & $O(TD)$/step & No \\
RNN/LSTM     & $O(TD^2)$     & $O(D)$     & $O(D^2)$/step & No \\
Mamba/SSM    & $O(TD)$       & $O(D)$     & $O(D)$/step  & Yes \\
FM (ours)    & $O(TD)$       & $O(D)$     & $O(D)$/step  & No \\
\bottomrule
\end{tabular}
\caption{Asymptotic complexity comparison. $T$: sequence length, $D$: hidden dimension, $L$: number of layers.}
\end{table}

FM matches Mamba's asymptotic complexity without requiring custom CUDA kernels, and exceeds transformer efficiency at long sequence lengths.

\subsection{Limitations}

\textbf{Exact recall}: the field is a lossy compression of the sequence. FM cannot retrieve the exact identity of a specific early token from the accumulated field, as that token's contribution is mixed with all subsequent contributions. This makes FM unsuitable for tasks requiring verbatim recall of specific positions.

\textbf{Long-range resolution}: as the sequence grows, each token's contribution becomes one of $T$ terms in the sum. While position orthogonality maintains separability in principle, the practical resolution of the decoder over very long sequences is an empirical question not fully explored in this work.

\textbf{Commutative baseline}: summation is commutative — $a + b = b + a$. Order is encoded only through the position modulation $p_t$. If the position encoding does not provide sufficient separation, the model may fail to distinguish different orderings. The multiplicative fusion and orthogonal position encoding are designed to address this, but formal guarantees on ordering sensitivity are left to future work.

\section{Future Work}

\textbf{Language modeling}: the FM architecture is domain-agnostic. The DNA encoding is specific to MIDI, but the field accumulation and position encoding are applicable to any discrete sequence. Replacing the domain-specific DNA with learned embeddings (or structured embeddings for other domains) and evaluating on standard language modeling benchmarks (PTB, WikiText-103, The Pile) is the immediate next step.

\textbf{Scaling}: FM has been validated at 13.78M parameters. Whether the architecture scales favorably — whether loss continues to decrease with more parameters and data — is an open empirical question. The linear training complexity suggests that scaling to longer sequences is less costly than for transformers.

\textbf{Hierarchical FM}: the architecture naturally extends to hierarchical processing. A fine-grained FM processes individual notes; a coarse FM processes phrase-level summaries; the coarse field informs the fine-grained decoding. This mirrors the multi-timescale structure of music and language.

\textbf{Interpretability}: the field $f_t$ is a single vector encoding the entire context up to position $t$. The trajectory of $f_t$ through $\mathbb{R}^D$ as a sequence is processed is the model's full computational history — a path in high-dimensional space. This opens geometric interpretability methods not available for transformers: attractor analysis, trajectory clustering, perturbation studies.

\textbf{Continual learning}: FM's additive field update is naturally compatible with continual learning. New sequences extend the field without rewriting it. Whether this translates to reduced catastrophic forgetting compared to transformer fine-tuning is an empirical question worth investigating.

\textbf{Edge deployment}: at inference, the model state is 8KB (4096 dimensions, float16). This is smaller than most L1 caches. FM's inference characteristics — fixed memory, constant compute, no attention, no KV cache — make it a candidate for deployment on severely memory-constrained hardware.

\section{Conclusion}

We have presented the Field Machine, a sequence architecture that achieves full training parallelism and $O(1)$ inference without recurrence, attention, or custom hardware operations. The key contributions are:

\begin{enumerate}
    \item \textbf{Structured token DNA}: domain geometry encoded directly into token representations before training, eliminating the need to learn geometric relationships from co-occurrence statistics.

    \item \textbf{Multiplicative position fusion}: position and identity fused into a single inseparable object, enabling the summation to preserve ordering information.

    \item \textbf{Cumulative field accumulation}: the entire sequence processed in one parallel \texttt{cumsum} operation, with no serial dependency anywhere in the forward pass.

    \item \textbf{$O(1)$ inference}: state is a single fixed-size vector, updated by one addition per token, independent of sequence length.
\end{enumerate}

These four primitives together constitute a new sequence modeling paradigm. The architecture was validated on symbolic music generation, reaching loss 3.09 in 4 minutes 43 seconds on a consumer GPU, crossing below the random baseline within two epochs.

The core insight is that a sequence is not inherently a chain. It is an ordered set. Order is a label. When that label is encoded directly into each element's representation, the chain dissolves — and with it, the serial bottleneck that has constrained sequence modeling since its inception.

\bibliographystyle{plain}
\begin{thebibliography}{99}

\bibitem{vaswani2017attention}
Vaswani, A., Shazeer, N., Parmar, N., Uszkoreit, J., Jones, L., Gomez, A. N., Kaiser, L., \& Polosukhin, I. (2017).
\textit{Attention is all you need}.
Advances in Neural Information Processing Systems, 30.

\bibitem{hochreiter1997long}
Hochreiter, S., \& Schmidhuber, J. (1997).
\textit{Long short-term memory}.
Neural Computation, 9(8), 1735--1780.

\bibitem{cho2014learning}
Cho, K., Van Merriënboer, B., Gulcehre, C., Bahdanau, D., Bougares, F., Schwenk, H., \& Bengio, Y. (2014).
\textit{Learning phrase representations using RNN encoder-decoder for statistical machine translation}.
arXiv preprint arXiv:1406.1078.

\bibitem{gu2023mamba}
Gu, A., \& Dao, T. (2023).
\textit{Mamba: Linear-time sequence modeling with selective state spaces}.
arXiv preprint arXiv:2312.00752.

\bibitem{peng2023rwkv}
Peng, B., Alcaide, E., Anthony, Q., Albalak, A., Arcadinho, S., Cao, H., ... \& Zhao, R. (2023).
\textit{RWKV: Reinventing RNNs for the transformer era}.
arXiv preprint arXiv:2305.13048.

\bibitem{dao2022flashattention}
Dao, T., Fu, D. Y., Ermon, S., Rudra, A., \& Ré, C. (2022).
\textit{FlashAttention: Fast and memory-efficient exact attention with IO-awareness}.
Advances in Neural Information Processing Systems, 35.

\bibitem{katharopoulos2020transformers}
Katharopoulos, A., Vyas, A., Pappas, N., \& Fleuret, F. (2020).
\textit{Transformers are RNNs: Fast autoregressive transformers with linear attention}.
International Conference on Machine Learning.

\bibitem{child2019generating}
Child, R., Gray, S., Radford, A., \& Sutskever, I. (2019).
\textit{Generating long sequences with sparse transformers}.
arXiv preprint arXiv:1904.10509.

\bibitem{su2022roformer}
Su, J., Lu, Y., Pan, S., Murtadha, A., Wen, B., \& Liu, Y. (2022).
\textit{RoFormer: Enhanced transformer with rotary position embedding}.
arXiv preprint arXiv:2104.09864.

\bibitem{hamilton2017embedding}
Hamilton, W. L., Leskovec, J., \& Jurafsky, D. (2017).
\textit{Embedding methods for link prediction}.
arXiv preprint arXiv:1709.05584.

\bibitem{raissi2019physics}
Raissi, M., Perdikaris, P., \& Karniadakis, G. E. (2019).
\textit{Physics-informed neural networks: A deep learning framework for solving forward and inverse problems involving nonlinear partial differential equations}.
Journal of Computational Physics, 378, 686--707.

\bibitem{gsm2026}
Grace. (2026).
\textit{GSM: Geometric State Machine — a novel sequence architecture built on fixed-manifold transformation algebras}.
Unpublished manuscript.

\end{thebibliography}

\end{document}
